\section{Introduction}

\subsection{Background}

Charts and graphs are among the most common methods for presenting quantitative information
in scientific publications, business reports, and data dashboards. A typical research paper
contains between two and ten chart images, each encoding numerical relationships that are
critical for understanding experimental results. Despite advances in machine learning and
natural language processing, the task of automatically extracting structured numerical data
from chart images remains a significant challenge.

Current approaches to automatic chart understanding fall into two categories. The first
category uses \textbf{end-to-end multimodal models} such as DePlot \cite{liu2023deplot} and
MatCha \cite{lee2023matcha}, which feed chart images directly into large vision-language
models that predict data tables or answer questions. While these approaches benefit from the
generalization capabilities of large models, they suffer from a fundamental limitation:
the model \emph{estimates} numerical values from pixel patterns rather than \emph{measuring}
them from geometric structures. This leads to hallucination errors where the model fabricates
plausible but incorrect data points.

The second category relies on \textbf{rule-based geometric analysis}, such as ReVision
\cite{savva2011revision}, which uses Hough transforms and template matching to extract
values from specific chart types. These methods are precise when applicable but lack
generalizability across chart types and fail on noisy or complex layouts.

Neither paradigm alone is sufficient. End-to-end models lack geometric precision, while
rule-based systems lack the perceptual flexibility of neural networks and the semantic
understanding of language models. This gap motivates a \textbf{hybrid neuro-symbolic
approach} that combines the strengths of both paradigms.

Furthermore, existing chart understanding systems are designed almost exclusively for
English-language charts. In the Vietnamese context, academic publications, government
reports, and business dashboards increasingly use charts with Vietnamese labels, axis
titles, and legends -- including diacritical marks (e.g., ``Doanh thu'',
``Tỷ lệ phần trăm'',
``Biểu đồ cột''). No published
system addresses Vietnamese chart understanding with structured data extraction and
Vietnamese-language question answering. This project targets this gap by designing the
pipeline as a \textbf{language-extensible} system: the core geometric and neural
components are language-agnostic, while OCR configuration, prompt templates, and SLM
fine-tuning data serve as localization layers that can be switched to Vietnamese without
modifying core logic.

\subsection{Objective}

This project aims to design and implement \textbf{Geo-SLM}, a hybrid AI system for
extracting structured numerical data from chart images. The system combines three
complementary paradigms:

\begin{enumerate}
    \item \textbf{Deep Learning for perception}: YOLOv8 \cite{jocher2023yolov8} for chart
          detection and ResNet-18 \cite{he2016resnet} for chart type classification.
    \item \textbf{Classical geometry for precision}: RANSAC-based axis calibration
          \cite{fischler1981ransac}, Ramer-Douglas-Peucker polyline simplification
          \cite{douglas1973rdp, ramer1972rdp}, and skeleton-based structural analysis
          \cite{lee1994skeleton} using OpenCV \cite{bradski2000opencv}.
    \item \textbf{Small Language Model (SLM) for reasoning}: A fine-tuned Qwen-2.5-1.5B
          \cite{yang2024qwen2} using QLoRA \cite{dettmers2024qlora} for OCR correction,
          value validation, and natural language description generation.
\end{enumerate}

The specific objectives are:
\begin{itemize}
    \item Achieve chart detection mAP@0.5 $> 85\%$ and classification accuracy $> 90\%$.
    \item Extract structured data (axis values, data series, labels) from 8 chart types
          with extraction confidence $> 90\%$.
    \item Build a large-scale training dataset ($> 200{,}000$ samples) from academic papers
          for SLM fine-tuning.
    \item Demonstrate that the hybrid geometric-plus-AI approach outperforms pure end-to-end
          methods in numerical accuracy while running locally on consumer hardware
          (NVIDIA RTX 3060, 6~GB VRAM).
    \item Enable \textbf{Vietnamese-language chart question answering}: the system must
          process charts containing Vietnamese text and support interactive Q\&A in
          Vietnamese, leveraging the extensible pipeline architecture.
\end{itemize}

\subsection{Problem Statement}

Given a document (PDF or image) containing one or more chart images, the system must:

\begin{itemize}
    \item \textbf{Input}: A PDF document or image file (PNG, JPG) containing charts.
    \item \textbf{Output}: For each detected chart, a structured JSON object containing:
          chart type, title, axis labels, data series (named lists of $(x, y)$ points),
          natural language description, and confidence scores with full traceability
          to source coordinates. The system must also support Vietnamese-language
          question answering over extracted chart data.
    \item \textbf{Chart types}: 8 categories -- line, bar, scatter, pie, histogram, box,
          heatmap, and area.
    \item \textbf{Performance targets}: Detection mAP@0.5 $> 85\%$; Classification accuracy
          $> 90\%$; End-to-end extraction confidence $> 80\%$.
    \item \textbf{Deployment constraint}: Must run on a single consumer GPU (RTX 3060,
          6~GB VRAM) without requiring cloud API access for core inference.
\end{itemize}

\subsection{Scope and Limitations}

\textbf{In scope:}
\begin{itemize}
    \item Charts from academic papers (arXiv cs.CV, cs.LG, cs.AI) and Vietnamese-language
          documents.
    \item 8 chart types: line, bar, scatter, pie, histogram, box, heatmap, area.
    \item PDF and raster image inputs (PNG, JPG, TIFF, BMP).
    \item Structured data extraction with JSON output.
    \item Local inference using fine-tuned SLM (1.5B parameters).
    \item Vietnamese-language OCR, chart Q\&A, and description generation.
\end{itemize}

\textbf{Out of scope:}
\begin{itemize}
    \item Business dashboards, presentation slides, and infographics.
    \item 3D charts, Sankey diagrams, and multi-panel composite figures.
    \item Real-time streaming or interactive chart manipulation.
    \item DOCX and HTML input formats (planned for future work).
\end{itemize}

\textbf{Known limitations:}
\begin{itemize}
    \item The core training data is sourced from English academic publications; Vietnamese
          chart data is generated through the localization pipeline (see \Cref{sec:viet_loc}).
    \item Axis calibration requires at least two readable numeric tick labels; charts
          without visible axis numbers produce no calibrated values.
    \item The SLM fine-tuning is ongoing at the time of writing; final SLM evaluation
          results (both English and Vietnamese) will be supplemented upon completion.
    \item Vietnamese OCR accuracy depends on diacritical mark recognition quality,
          which varies across OCR engines and font styles.
\end{itemize}

